Let a moving coordinate system $K$ undergo general motion (translation plus rotation) relative to a stationary system $k$. Let $B_t: K \to k$ be the linear operator describing the rotation of $K$ relative to $k$, let $r(t) \in k$ be the position of the origin of $K$, and let $\omega$ be the angular velocity vector of $K$ in $k$. Let $Q(t) \in K$ be the radius vector of a moving point in the moving coordinate system, and $q(t) = B_t Q(t) + r(t) \in k$ be the radius vector in the stationary system. Then the absolute velocity decomposes as
$$
v = v' + v_n + v_0,
$$
where
$$
v = \dot{q} \in k \text{ is the absolute velocity},
$$
$$
v' = B \dot{Q} \in k \text{ is the relative velocity},
$$
$$
v_n = \dot{B} Q = [\omega, q - r] \in k \text{ is the transferred velocity of rotation},
$$
$$
v_0 = \dot{r} \in k \text{ is the velocity of motion of the moving coordinate system}.
$$
